\documentclass[12pt]{article}

% Packages
\usepackage[margin=2.5cm]{geometry}
\usepackage{amsmath,amssymb,amsthm}
\usepackage{bm}
\usepackage{enumitem}
\usepackage{booktabs}
\usepackage{graphicx}
\usepackage[hidelinks]{hyperref}

% Theorem environments
\newtheorem{proposition}{Proposition}
\newtheorem{theorem}{Theorem}
\newtheorem{corollary}{Corollary}
\newtheorem{remark}{Remark}

% No indentation
\setlength{\parindent}{0pt}
\setlength{\parskip}{6pt}

\begin{document}

\title{Convergence of the Universal Route\\
to Classical Isotherm Models\\[6pt]
\large (Version 4 --- Limiting-case proofs)}
\author{}
\date{}
\maketitle

% ======================================================================
\section{Introduction}
% ======================================================================

In the preceding versions (v2--v3) we established a model-agnostic
protocol for extracting $Q_{\max}$ and
$K_{\mathrm{aff}}\equiv K_H/Q_{\max}$ from adsorption data, based on
the heterogeneous Langmuir integral
\begin{equation}
  q(C) = K_p\,C + \int_{0}^{\infty} \rho(K)\,\frac{KC}{1+KC}\,dK,
  \qquad \rho(K)\ge 0,
  \label{eq:master}
\end{equation}
with invariants
\begin{equation}
  Q_{\max} = \int_0^\infty \rho(K)\,dK, \qquad
  K_H = \int_0^\infty K\,\rho(K)\,dK, \qquad
  K_{\mathrm{aff}} = \frac{K_H}{Q_{\max}}.
  \label{eq:invariants}
\end{equation}

The purpose of this document is to \emph{prove} that every classical
adsorption isotherm that satisfies assumptions A1--A4 is a special
case of Eq.~\eqref{eq:master}, and that the universal invariants
$Q_{\max}$, $K_H$, and $K_{\mathrm{aff}}$ reduce exactly to the
known parameters of each model.  We also characterise the degenerate
cases (Freundlich, Redlich--Peterson with $\beta<1$) where one or more
invariants diverge, and prove a kernel-independence theorem that makes
the invariants robust to the choice of local isotherm mechanism.

% ======================================================================
\section{Framework recap}
\label{sec:framework}
% ======================================================================

\textbf{Assumptions} (from v2--v3):
\begin{itemize}[nosep]
  \item[A1.] Independent-site superposition (heterogeneous Langmuir
        class, no cooperativity).
  \item[A2.] Finite saturating capacity:
        $Q_{\max}=\int_0^\infty\rho(K)\,dK<\infty$.
  \item[A3.] Finite first moment:
        $K_H=\int_0^\infty K\rho(K)\,dK<\infty$.
  \item[A4.] Sufficient experimental coverage of the Henry and plateau
        regimes.
\end{itemize}

The \emph{normalised} affinity distribution is
$f(K)\equiv\rho(K)/Q_{\max}$, with $\int_0^\infty f(K)\,dK=1$.
The effective affinity is the first moment of $f$:
\begin{equation}
  K_{\mathrm{aff}} = \int_0^\infty K\,f(K)\,dK = \langle K\rangle_f.
  \label{eq:Kaff_moment}
\end{equation}

% ======================================================================
\section{Exact convergence to classical models}
\label{sec:convergence}
% ======================================================================

% ------ 3.1 Langmuir ------
\subsection{Langmuir isotherm}

\begin{proposition}[Langmuir]
\label{prop:langmuir}
Let $\rho(K)=Q_m\,\delta(K-K_0)$ with $Q_m>0$ and $K_0>0$.  Then:
\begin{enumerate}[nosep]
  \item $q(C) = \displaystyle\frac{Q_m K_0 C}{1+K_0 C}$
        \quad (single-site Langmuir),
  \item $Q_{\max}=Q_m$, \quad $K_H=Q_m K_0$, \quad
        $K_{\mathrm{aff}}=K_0$.
\end{enumerate}
\end{proposition}

\begin{proof}
Substituting $\rho(K)=Q_m\,\delta(K-K_0)$ into
Eq.~\eqref{eq:master} (with $K_p=0$):
\[
  q(C)=\int_0^\infty Q_m\,\delta(K-K_0)\,\frac{KC}{1+KC}\,dK
      = Q_m\,\frac{K_0 C}{1+K_0 C}.
\]
The invariants follow directly:
$Q_{\max}=\int Q_m\,\delta(K-K_0)\,dK=Q_m$,\;
$K_H=\int K\,Q_m\,\delta(K-K_0)\,dK=Q_m K_0$,\;
$K_{\mathrm{aff}}=K_0$.
All assumptions A1--A4 are satisfied.
\end{proof}

% ------ 3.2 Multi-site Langmuir ------
\subsection{Multi-site Langmuir isotherm}

\begin{proposition}[Multi-site Langmuir]
\label{prop:multilang}
Let $\rho(K)=\sum_{j=1}^{J}Q_j\,\delta(K-K_j)$ with
$Q_j>0$, $K_j>0$ distinct.  Then:
\begin{enumerate}[nosep]
  \item $q(C)=\displaystyle\sum_{j=1}^{J}
        \frac{Q_j K_j C}{1+K_j C}$
        \quad ($J$-site Langmuir),
  \item $Q_{\max}=\displaystyle\sum_{j=1}^{J}Q_j$, \quad
        $K_H=\displaystyle\sum_{j=1}^{J}Q_j K_j$, \quad
        $K_{\mathrm{aff}}=\displaystyle
        \frac{\sum_j Q_j K_j}{\sum_j Q_j}$.
\end{enumerate}
\end{proposition}

\begin{proof}
Linearity of the integral and the sifting property of the
Dirac delta give (1) immediately.  The invariants are finite sums
of positive terms, hence finite.  $K_{\mathrm{aff}}$ is the
capacity-weighted arithmetic mean of the individual affinities.
\end{proof}

\begin{corollary}[Bounds on $K_{\mathrm{aff}}$]
\label{cor:bounds}
For any discrete mixture,
$\min_j K_j \le K_{\mathrm{aff}} \le \max_j K_j$.
\end{corollary}

\begin{proof}
This follows from $K_{\mathrm{aff}}$ being a convex combination
of the $K_j$ with positive weights $Q_j/Q_{\max}$.
\end{proof}

% ------ 3.3 Sips ------
\subsection{Sips (Langmuir--Freundlich) isotherm}

\begin{proposition}[Sips]
\label{prop:sips}
Let the affinity distribution be the Sips
distribution~\cite{Sips1948}:
\begin{equation}
  f(K) = \frac{\sin(m\pi)}{\pi}\,
  \frac{(K/K_s)^{m}}
       {K\bigl[(K/K_s)^{2m}+2\cos(m\pi)(K/K_s)^{m}+1\bigr]},
  \qquad 0<m\le 1,\; K_s>0.
  \label{eq:sips_dist}
\end{equation}
Then:
\begin{enumerate}[nosep]
  \item $q(C) = \displaystyle\frac{Q_{\max}(K_s C)^m}
        {1+(K_s C)^m}$
        \quad (Sips isotherm).
  \item $Q_{\max}$ is finite and correctly recovered:
        $\lim_{C\to\infty}q(C)=Q_{\max}$.
  \item For $m=1$: the Sips distribution degenerates to
        $\delta(K-K_s)$, so $K_H=Q_{\max}K_s$ and
        $K_{\mathrm{aff}}=K_s$ (Langmuir limit).
  \item For $0<m<1$: $K_H=+\infty$ (first moment diverges).
        Assumption A3 is violated.
\end{enumerate}
\end{proposition}

\begin{proof}
\textit{Part (1):} The isotherm follows from the Stieltjes-transform
inversion by Sips~\cite{Sips1948}; see also
Jaroniec and Madey~\cite{Jaroniec1988}, \S3.2.

\textit{Part (2):} Since $(K_s C)^m/(1+(K_s C)^m)\to 1$ as
$C\to\infty$, the limit gives $Q_{\max}$.

\textit{Part (3):} As $m\to 1$, $\sin(m\pi)\to 0$ and the
distribution narrows; by L'H\^opital's rule the Sips distribution
converges to $\delta(K-K_s)$ in the distributional
sense~\cite{Jaroniec1988}.

\textit{Part (4):} For $0<m<1$, differentiate the isotherm:
\[
  \frac{dq}{dC}\bigg|_{C\to 0}
  = Q_{\max}\,m\,K_s^m\,\lim_{C\to 0^+}C^{m-1}
  = +\infty,
\]
since $m-1<0$.  Equivalently, the tail of $f(K)$ decays as
$f(K)\sim K^{-m-1}$ for $K\to\infty$, so
$\int_0^\infty K f(K)\,dK$ diverges.
\end{proof}

\begin{remark}[Operational $K_{\mathrm{aff}}$ for Sips]
Although $K_H=\infty$ for $m<1$, the universal workflow produces a
\emph{finite} operational $\widehat{K}_{\mathrm{aff}}$ from data
spanning a finite concentration range $[C_{\min},C_{\max}]$.  This
estimate depends on the range and should be reported together with
the concentration window used.  The sensitivity diagnostic
(Section~5.3 of v3) flags this case: $K_{\mathrm{aff}}$ varies
substantially with the size of the Henry region $\mathcal{I}_0$.
\end{remark}

% ------ 3.4 Tóth ------
\subsection{T\'oth isotherm}

\begin{proposition}[T\'oth]
\label{prop:toth}
The T\'oth isotherm~\cite{Toth1971}
\begin{equation}
  q(C) = \frac{Q_{\max}\,K_T\,C}{\bigl(1+(K_T C)^t\bigr)^{1/t}},
  \qquad 0<t\le 1,\; K_T>0,
  \label{eq:toth}
\end{equation}
satisfies the heterogeneous Langmuir integral with a specific
$\rho(K)$ related to one-sided L\'evy stable
distributions~\cite{Jaroniec1975}.  Its invariants are:
\begin{enumerate}[nosep]
  \item $Q_{\max}=Q_{\max}$ (finite for all $t$).
  \item $K_H = Q_{\max}\,K_T$ (finite for all $0<t\le 1$).
  \item $K_{\mathrm{aff}} = K_T$.
  \item For $t=1$: reduces to the Langmuir isotherm.
\end{enumerate}
\end{proposition}

\begin{proof}
\textit{Part (1):}
\[
  \lim_{C\to\infty} \frac{Q_{\max}K_T C}{(1+(K_T C)^t)^{1/t}}
  = \lim_{C\to\infty} \frac{Q_{\max}K_T C}{(K_T C)^{t/t}}
    \cdot\frac{1}{(1+(K_TC)^{-t})^{1/t}}
  = Q_{\max}\cdot 1 = Q_{\max}.
\]

\textit{Part (2):}  Differentiate at $C=0$.  Write
$q(C)=Q_{\max}K_T C\cdot g(C)$ where
$g(C)=(1+(K_T C)^t)^{-1/t}$.  Then $g(0)=1$ and
\[
  K_H = \frac{dq}{dC}\bigg|_{C=0} = Q_{\max}K_T\cdot g(0) = Q_{\max}K_T.
\]

\textit{Part (3):} $K_{\mathrm{aff}}=K_H/Q_{\max}=K_T$.

\textit{Part (4):} For $t=1$:
$(1+K_T C)^{1/1}=1+K_T C$, giving the standard Langmuir.
\end{proof}

\begin{remark}[Heterogeneity parameter $t$]
The T\'oth parameter $t$ controls asymmetry of the energy
distribution.  For $t<1$ the distribution is broader with a longer
tail toward weak sites ($K\to 0$).  Unlike the Sips parameter $m$,
the T\'oth parameter does \emph{not} cause the first moment to
diverge; hence $K_{\mathrm{aff}}=K_T$ is well-defined for all
$t\in(0,1]$.  This makes the T\'oth model particularly compatible
with the universal route.
\end{remark}

% ------ 3.5 Temkin ------
\subsection{Temkin isotherm}

\begin{proposition}[Temkin]
\label{prop:temkin}
Let the affinity distribution be uniform in $\ln K$\,:
\begin{equation}
  \rho(K) = \frac{Q_{\max}}{K\ln(K_{\max}/K_{\min})},
  \qquad K_{\min}\le K \le K_{\max},
  \label{eq:temkin_dist}
\end{equation}
with $0<K_{\min}<K_{\max}<\infty$.  Then:
\begin{enumerate}[nosep]
  \item $q(C) = \displaystyle\frac{Q_{\max}}{\ln(K_{\max}/K_{\min})}
        \bigl[\ln(1+K_{\max}C)-\ln(1+K_{\min}C)\bigr]$.
  \item $Q_{\max}=Q_{\max}$ (finite).
  \item $K_H = \displaystyle\frac{Q_{\max}(K_{\max}-K_{\min})}
        {\ln(K_{\max}/K_{\min})}$ (finite).
  \item $K_{\mathrm{aff}} = \displaystyle
        \frac{K_{\max}-K_{\min}}{\ln(K_{\max}/K_{\min})}$
        \quad (logarithmic mean of $K_{\min}$ and $K_{\max}$).
  \item In the intermediate range $K_{\min}C\ll 1\ll K_{\max}C$:
        $q\approx b\,\ln(K_{\max}C)$ with
        $b\equiv Q_{\max}/\ln(K_{\max}/K_{\min})$
        (classical Temkin form).
\end{enumerate}
\end{proposition}

\begin{proof}
\textit{Part (1):}
\begin{align*}
  q(C) &= \int_{K_{\min}}^{K_{\max}}
    \frac{Q_{\max}}{K\ln(K_{\max}/K_{\min})}
    \,\frac{KC}{1+KC}\,dK
  = \frac{Q_{\max}}{\ln(K_{\max}/K_{\min})}
    \int_{K_{\min}}^{K_{\max}}
    \frac{C}{1+KC}\,dK \\
  &= \frac{Q_{\max}}{\ln(K_{\max}/K_{\min})}
    \bigl[\ln(1+KC)\bigr]_{K_{\min}}^{K_{\max}}
  = \frac{Q_{\max}}{\ln(K_{\max}/K_{\min})}
    \bigl[\ln(1+K_{\max}C)-\ln(1+K_{\min}C)\bigr].
\end{align*}

\textit{Part (2):} As $C\to\infty$:
$\ln(1+K_{\max}C)-\ln(1+K_{\min}C)\to
\ln(K_{\max}C)-\ln(K_{\min}C)=\ln(K_{\max}/K_{\min})$,
so $q\to Q_{\max}$.

\textit{Part (3):}
\begin{align*}
  K_H &= \int_{K_{\min}}^{K_{\max}}
    K\cdot\frac{Q_{\max}}{K\ln(K_{\max}/K_{\min})}\,dK
  = \frac{Q_{\max}}{\ln(K_{\max}/K_{\min})}
    \int_{K_{\min}}^{K_{\max}} dK
  = \frac{Q_{\max}(K_{\max}-K_{\min})}
        {\ln(K_{\max}/K_{\min})}.
\end{align*}

\textit{Part (4):} Division of (3) by (2).

\textit{Part (5):} When $K_{\min}C\ll 1$,
$\ln(1+K_{\min}C)\approx K_{\min}C\approx 0$.
When $K_{\max}C\gg 1$,
$\ln(1+K_{\max}C)\approx \ln(K_{\max}C)$.
Hence $q\approx b\ln(K_{\max}C)$, which is the classical
Temkin equation $q=b\ln(AC)$ with $A=K_{\max}$.
\end{proof}

\begin{remark}[Logarithmic mean affinity]
The result $K_{\mathrm{aff}}=(K_{\max}-K_{\min})/
\ln(K_{\max}/K_{\min})$ is the \emph{logarithmic mean}
of $K_{\min}$ and $K_{\max}$, satisfying
$K_{\min}<K_{\mathrm{aff}}<K_{\max}$.
As $K_{\min}\to K_{\max}$ (homogeneous surface),
$K_{\mathrm{aff}}\to K_{\max}$ and the isotherm reduces to a
single-site Langmuir (Proposition~\ref{prop:langmuir}).
\end{remark}

% ------ 3.6 Dual-mode ------
\subsection{Dual-mode (partition + Langmuir) isotherm}

\begin{proposition}[Dual-mode]
\label{prop:dualmode}
Let $q(C)=K_p C + Q_m K_0 C/(1+K_0 C)$ with $K_p>0$, $Q_m>0$,
$K_0>0$.  Then:
\begin{enumerate}[nosep]
  \item The saturating component has
        $\rho(K)=Q_m\,\delta(K-K_0)$,
        and the linear component is $K_p C$.
  \item $Q_{\max}=Q_m$ (saturating capacity only).
  \item $K_H=Q_m K_0$ (Henry constant of the saturating part).
  \item $K_{\mathrm{aff}}=K_0$.
  \item The total initial slope is $K_p + K_H = K_p + Q_m K_0$;
        the universal workflow separates $K_p$ from $K_H$ by
        detecting the persistent linear trend at high~$C$.
\end{enumerate}
\end{proposition}

\begin{proof}
Parts (1)--(4) follow from Proposition~\ref{prop:langmuir}
applied to the saturating component.  Part (5) is the operational
strategy: at high $C$, $q\approx K_p C + Q_m$ (linear with
intercept $Q_m$ and slope $K_p$), so $K_p$ is identifiable from
the high-$C$ slope.  After subtracting $K_p C$, the residual
is a standard Langmuir.
\end{proof}

% ======================================================================
\section{Universal low-concentration limit}
\label{sec:henry}
% ======================================================================

\begin{proposition}[Henry law as universal limit]
\label{prop:henry}
For any distribution $\rho(K)\ge 0$ satisfying A2 and A3 (with
$K_p\ge 0$):
\begin{equation}
  q(C) = (K_p + K_H)\,C + O(C^2) \qquad\text{as } C\to 0.
  \label{eq:henry_universal}
\end{equation}
\end{proposition}

\begin{proof}
The Langmuir kernel admits the Taylor expansion
$KC/(1+KC) = KC - K^2 C^2 + O(C^3)$.  Under assumption A3,
$\int_0^\infty K\rho(K)\,dK = K_H<\infty$, so by dominated
convergence we may interchange the integral and the limit:
\begin{align*}
  q(C) &= K_p C + \int_0^\infty \rho(K)\bigl[KC - K^2 C^2
          + O(C^3)\bigr]\,dK \\
       &= K_p C + K_H\,C - \left(\int_0^\infty K^2\rho(K)\,dK\right)C^2
          + O(C^3).
\end{align*}
If additionally the second moment
$\int K^2\rho(K)\,dK<\infty$, the $O(C^2)$ coefficient is
well-defined and equals $-Q_{\max}\langle K^2\rangle_f$.
\end{proof}

\begin{corollary}
The Henry constant $K_H$ is a model-independent invariant: it equals
$Q_{\max}\langle K\rangle_f$ regardless of the shape of $f(K)$,
as long as the first moment exists.
\end{corollary}

% ======================================================================
\section{Kernel-independence theorem}
\label{sec:kernel}
% ======================================================================

The invariants $Q_{\max}$ and $K_H$ (hence $K_{\mathrm{aff}}$) were
derived using the Langmuir kernel $\varphi(C,K)=KC/(1+KC)$.
The following theorem shows they are \emph{independent} of the
specific local isotherm mechanism.

\begin{theorem}[Kernel independence of asymptotic invariants]
\label{thm:kernel}
Let $\varphi(C,K)$ be \emph{any} local isotherm kernel satisfying:
\begin{enumerate}[nosep]
  \item[(K1)] $\varphi(0,K)=0$ for all $K>0$
        \quad (no uptake at zero concentration),
  \item[(K2)] $\displaystyle\frac{\partial\varphi}{\partial C}
        \bigg|_{C=0}=K$ for all $K>0$
        \quad (Henry-law onset with constant $K$),
  \item[(K3)] $\displaystyle\lim_{C\to\infty}\varphi(C,K)=1$
        for all $K>0$
        \quad (saturation of each site class),
  \item[(K4)] $0\le\varphi(C,K)\le 1$ for all $C,K\ge 0$
        \quad (bounded).
\end{enumerate}
Consider the generalised heterogeneous isotherm
\begin{equation}
  q(C)=\int_0^\infty\rho(K)\,\varphi(C,K)\,dK.
  \label{eq:general_kernel}
\end{equation}
Then, under assumptions A2 and A3:
\begin{equation}
  \lim_{C\to\infty}q(C)=Q_{\max},\qquad
  \frac{dq}{dC}\bigg|_{C=0}=K_H,\qquad
  K_{\mathrm{aff}}=\frac{K_H}{Q_{\max}}.
  \label{eq:kernel_indep}
\end{equation}
\end{theorem}

\begin{proof}
\textit{Saturation limit:}
By (K3), $\varphi(C,K)\to 1$ for each $K>0$ as $C\to\infty$.
By (K4), $|\rho(K)\varphi(C,K)|\le\rho(K)$, which is integrable by
A2.  By the dominated convergence theorem:
\[
  \lim_{C\to\infty}q(C) = \int_0^\infty \rho(K)\cdot 1\,dK
  = Q_{\max}.
\]

\textit{Henry slope:}
By (K1) and (K2), $\varphi(C,K)/C\to K$ as $C\to 0$.
By (K4), $\rho(K)\varphi(C,K)/C\le\rho(K)\cdot K$ (for $C$
small enough), which is integrable by A3.  By dominated convergence:
\[
  \lim_{C\to 0}\frac{q(C)}{C}
  = \int_0^\infty \rho(K)\cdot K\,dK = K_H.
\]

\textit{Effective affinity:}
$K_{\mathrm{aff}}=K_H/Q_{\max}$ by definition.
\end{proof}

\begin{remark}[Examples of valid kernels]
In addition to the Langmuir kernel $\varphi_L=KC/(1+KC)$,
Theorem~\ref{thm:kernel} applies to:
\begin{itemize}[nosep]
  \item \textbf{Jovanovic kernel:}
        $\varphi_J(C,K)=1-e^{-KC}$.
        Check: $\varphi_J(0,K)=0$, $\varphi_J'(0)=K$,
        $\varphi_J(\infty)=1$.
  \item \textbf{Volmer kernel:}
        Any kernel derived from a 2-D equation of state for the
        adsorbed phase, provided it satisfies (K1)--(K4).
\end{itemize}
The NNLS distribution recovery \emph{does} depend on the kernel
choice, but the asymptotic invariants $Q_{\max}$, $K_H$, and
$K_{\mathrm{aff}}$ do not.
\end{remark}

% ======================================================================
\section{Degenerate cases: divergent invariants}
\label{sec:degenerate}
% ======================================================================

% ------ 6.1 Freundlich ------
\subsection{Pure Freundlich isotherm}

\begin{proposition}[Freundlich --- doubly degenerate]
\label{prop:freundlich}
The Freundlich isotherm $q(C)=K_F C^{n}$ with $0<n<1$ cannot be
represented by the heterogeneous Langmuir integral with finite
$Q_{\max}$ and $K_H$.  Specifically:
\begin{enumerate}[nosep]
  \item $Q_{\max}=\lim_{C\to\infty}q(C)=+\infty$
        \quad (A2 violated).
  \item $K_H=\lim_{C\to 0}dq/dC=+\infty$
        \quad (A3 violated).
  \item The underlying distribution $f(K)\propto K^{n-2}$ has both
        zeroth and first moments infinite.
\end{enumerate}
\end{proposition}

\begin{proof}
\textit{Part (1):} $\lim_{C\to\infty}K_F C^n=+\infty$ for $n>0$.

\textit{Part (2):}
$dq/dC=K_F n C^{n-1}\to+\infty$ as $C\to 0^+$ for $n<1$.

\textit{Part (3):}
The Freundlich isotherm is the low-coverage limit of the Sips
isotherm with $m=n$ and $Q_{\max}\to\infty$ (see
Proposition~\ref{prop:sips}).
The Sips distribution $f(K)$ has a tail $\sim K^{-m-1}$,
whose zeroth moment over $[0,\infty)$ diverges when $Q_{\max}$ is
not finite.
\end{proof}

\begin{remark}[Operational treatment of Freundlich data]
When the universal workflow is applied to data that follow a
Freundlich isotherm over a finite range $[C_{\min},C_{\max}]$:
\begin{itemize}[nosep]
  \item $\widehat{Q}_{\max}$ will be approximately $K_F C_{\max}^n$
        (the value at the highest measured concentration).
  \item $\widehat{K}_H$ will depend strongly on $C_{\min}$.
  \item The sensitivity diagnostic will flag
        $K_{\mathrm{aff}}$ as range-dependent.
  \item These operational estimates are \emph{not} material constants
        but depend on the experimental window.
\end{itemize}
\end{remark}

% ------ 6.2 Redlich-Peterson ------
\subsection{Redlich--Peterson isotherm}

\begin{proposition}[Redlich--Peterson]
\label{prop:rp}
The Redlich--Peterson isotherm~\cite{Redlich1959}
\begin{equation}
  q(C) = \frac{K_R\,C}{1+a_R\,C^{\beta}},
  \qquad 0<\beta\le 1,
  \label{eq:rp}
\end{equation}
has the following properties:
\begin{enumerate}[nosep]
  \item For $\beta=1$: reduces to Langmuir with
        $Q_{\max}=K_R/a_R$ and $K=a_R$.
  \item For $\beta<1$:
        $Q_{\max}=\lim_{C\to\infty}q(C)=+\infty$ (A2 violated),
        but $K_H=K_R$ is finite.
  \item The Redlich--Peterson isotherm with $\beta<1$ does not arise
        from the heterogeneous Langmuir integral with finite
        $Q_{\max}$.
\end{enumerate}
\end{proposition}

\begin{proof}
\textit{Part (1):} Set $\beta=1$:
$q=K_R C/(1+a_R C)=(K_R/a_R)\cdot a_R C/(1+a_R C)$, which is
Langmuir.

\textit{Part (2):}
$\lim_{C\to\infty}K_R C/(a_R C^\beta)
=(K_R/a_R)C^{1-\beta}\to\infty$ for $\beta<1$.
At $C\to 0$: $q\approx K_R C/(1+0)=K_R C$, so $K_H=K_R$.

\textit{Part (3):}
Since $Q_{\max}=\infty$, the isotherm cannot be written as
Eq.~\eqref{eq:master} with $\int\rho(K)\,dK<\infty$.
\end{proof}

% ======================================================================
\section{Models excluded by assumption A1}
\label{sec:excluded}
% ======================================================================

The following models violate Assumption A1
(independent-site superposition) and are therefore \emph{outside}
the scope of the universal route:

\begin{enumerate}[nosep]
  \item \textbf{BET isotherm}~\cite{BET1938}:
        Multilayer adsorption.  Each adsorbed molecule provides a
        new site for the next layer, violating independence.
        $Q_{\max}$ in the BET sense refers to monolayer coverage,
        but the isotherm diverges as $C\to C_s$ (saturation
        concentration).

  \item \textbf{Hill isotherm}
        $q=Q_{\max}(KC)^{n_H}/(1+(KC)^{n_H})$ with $n_H>1$:
        Positive cooperativity produces a sigmoidal isotherm.
        At $C=0$, $dq/dC=0$ for $n_H>1$, meaning $K_H=0$.
        The inflection point violates the monotone-concave shape
        expected from heterogeneous Langmuir superposition.

  \item \textbf{Fowler--Guggenheim isotherm}:
        Lateral interactions between adsorbed molecules
        ($q$-dependent free energy) produce a non-Langmuir local
        isotherm.  For strong enough interactions, the isotherm
        can become S-shaped or exhibit hysteresis.
\end{enumerate}

\textbf{Diagnostic:}
If the data show a sigmoidal shape (inflection point), the isotherm
lies outside the heterogeneous Langmuir class, and the universal
route should not be applied.  A simple test: verify that $d^2 q/dC^2<0$
for all measured $C>0$.

% ======================================================================
\section{Summary: classification of isotherm models}
\label{sec:summary}
% ======================================================================

Table~\ref{tab:summary} classifies the main adsorption isotherm
models with respect to the universal route.

\begin{table}[htbp]
\centering
\caption{Convergence of the universal route to classical isotherms.
``$\checkmark$'' = finite and correctly recovered;
``$\infty$'' = divergent (assumption violated);
``---'' = not applicable.}
\label{tab:summary}
\small
\begin{tabular}{@{}lcccccl@{}}
\toprule
\textbf{Model} & \textbf{A1--A4} & $Q_{\max}$ & $K_H$
  & $K_{\mathrm{aff}}$ & $K_p$ & \textbf{Proposition} \\
\midrule
Langmuir
  & all & $Q_m$ & $Q_m K$ & $K$
  & $0$ & \ref{prop:langmuir} \\[2pt]
Multi-Langmuir ($J$-site)
  & all & $\sum Q_j$ & $\sum Q_j K_j$
  & $\frac{\sum Q_j K_j}{\sum Q_j}$
  & $0$ & \ref{prop:multilang} \\[2pt]
Sips ($m=1$)
  & all & $Q_m$ & $Q_m K_s$ & $K_s$
  & $0$ & \ref{prop:sips} \\[2pt]
Sips ($m<1$)
  & A3 viol. & $Q_m\;\checkmark$ & $\infty$
  & oper. & $0$ & \ref{prop:sips} \\[2pt]
T\'oth ($0<t\le 1$)
  & all & $Q_m$ & $Q_m K_T$ & $K_T$
  & $0$ & \ref{prop:toth} \\[2pt]
Temkin
  & all & $Q_m$ & $Q_m K_{\log}$
  & $K_{\log}$ & $0$ & \ref{prop:temkin} \\[2pt]
Dual-mode
  & all & $Q_m$ & $Q_m K$ & $K$
  & $K_p$ & \ref{prop:dualmode} \\[2pt]
Freundlich
  & A2,A3 viol. & $\infty$ & $\infty$
  & --- & --- & \ref{prop:freundlich} \\[2pt]
Redlich--Peterson ($\beta<1$)
  & A2 viol. & $\infty$ & $K_R\;\checkmark$
  & --- & $0$ & \ref{prop:rp} \\[2pt]
\midrule
BET & A1 viol. & --- & --- & --- & ---
  & Sect.~\ref{sec:excluded} \\
Hill ($n_H>1$)
  & A1 viol. & $Q_m$ & $0$ & $0$
  & --- & Sect.~\ref{sec:excluded} \\
Fowler--Guggenheim
  & A1 viol. & --- & --- & --- & ---
  & Sect.~\ref{sec:excluded} \\
\bottomrule
\end{tabular}

\medskip
$K_{\log}\equiv(K_{\max}-K_{\min})/\ln(K_{\max}/K_{\min})$
\quad (logarithmic mean).
\end{table}

\textbf{Key conclusions:}
\begin{enumerate}[nosep]
  \item Every isotherm within the heterogeneous Langmuir class
        (A1 satisfied) with finite $Q_{\max}$ and $K_H$
        (A2--A3 satisfied) yields exact convergence of the
        universal invariants to the model parameters.
  \item The T\'oth isotherm is the broadest single-equation model
        that satisfies all assumptions: it accommodates asymmetric
        heterogeneity while keeping $K_{\mathrm{aff}}=K_T$
        well-defined for all $t\in(0,1]$.
  \item The Sips isotherm with $m<1$ produces a finite
        $Q_{\max}$ but an infinite $K_H$; the universal workflow
        correctly recovers $Q_{\max}$ and flags the
        $K_{\mathrm{aff}}$ estimate as range-dependent.
  \item The Freundlich and Redlich--Peterson ($\beta<1$) isotherms
        have infinite $Q_{\max}$; the universal workflow gives
        operational estimates that are inherently range-dependent.
  \item Models violating A1 (BET, Hill, Fowler--Guggenheim) produce
        isotherms outside the concave-Langmuir class and are
        excluded.
\end{enumerate}

% ======================================================================
\section{Numerical verification}
\label{sec:numerical}
% ======================================================================

The companion MATLAB script
\texttt{universal\_route\_v4\_limits.m} verifies the
convergence results numerically.  For each model in
Table~\ref{tab:summary} (excluding the excluded models), synthetic
data with 5\% relative Gaussian noise are generated, the universal
workflow is applied, and the estimated invariants are compared with
the analytical values.

% ------ Figure 1 ------
\subsection{Isotherm reconstruction for all models}

Figure~\ref{fig:v4models} shows the reconstruction for seven
representative models.  In each panel, the grey curve is the true
generating model, black circles are the noisy data, the red dashed
curve is the effective Langmuir reconstruction using
$\widehat{Q}_{\max}$ and $\widehat{K}_{\mathrm{aff}}$, and the
green horizontal line marks $\widehat{Q}_{\max}$.

For the Langmuir, T\'oth, Temkin, and dual-mode models, the
reconstruction closely follows the true curve and the estimated
invariants match the analytical values within the noise level.
For the Sips model, $Q_{\max}$ is well estimated but
$K_{\mathrm{aff}}$ is operational (finite but range-dependent).
For the Freundlich model, both invariants are operational
approximations to the (infinite) true values.

\begin{figure}[htbp]
  \centering
  \includegraphics[width=\textwidth]{fig_v4_models.pdf}
  \caption{Isotherm reconstruction for seven models using the
  universal route.  Grey: true model; black circles: noisy data;
  red dashed: effective Langmuir reconstruction; green: estimated
  $\widehat{Q}_{\max}$.  The bottom-right panel shows the relative
  errors $|\widehat{Q}_{\max}-Q_{\max}^{\mathrm{true}}|
  /Q_{\max}^{\mathrm{true}}$ (blue) and
  $|\widehat{K}_{\mathrm{aff}}-K_{\mathrm{aff}}^{\mathrm{true}}|
  /K_{\mathrm{aff}}^{\mathrm{true}}$ (orange) for the models
  with finite invariants.}
  \label{fig:v4models}
\end{figure}

% ------ Figure 2 ------
\subsection{Convergence with decreasing noise}

Figure~\ref{fig:v4convergence} demonstrates that the estimation
errors in $Q_{\max}$ and $K_{\mathrm{aff}}$ decrease systematically
as the noise level is reduced, confirming consistency of the
universal workflow.

For models with well-defined invariants (Langmuir, T\'oth, Temkin,
dual-mode, double Langmuir), the relative error decreases roughly
linearly with the noise level on a log-log scale, indicating
$O(\sigma)$ convergence rate.
For the Sips model, $Q_{\max}$ converges but $K_{\mathrm{aff}}$
remains sensitive to the concentration range even at low noise,
confirming the first-moment divergence.

\begin{figure}[htbp]
  \centering
  \includegraphics[width=0.85\textwidth]{fig_v4_convergence.pdf}
  \caption{Relative error in $\widehat{Q}_{\max}$ (left) and
  $\widehat{K}_{\mathrm{aff}}$ (right) versus noise level
  for five models with finite invariants.  Each point is the
  median over 20 Monte Carlo realisations; error bars show
  the interquartile range.}
  \label{fig:v4convergence}
\end{figure}

% ======================================================================
\section*{References}
% ======================================================================

\begin{thebibliography}{99}

\bibitem{Sips1948}
R.~Sips, ``On the structure of a catalyst surface,''
\textit{J.\ Chem.\ Phys.}\ \textbf{16}, 490--495 (1948).

\bibitem{Jaroniec1988}
M.~Jaroniec and R.~Madey,
\textit{Physical Adsorption on Heterogeneous Solids},
Elsevier, Amsterdam (1988).

\bibitem{Rudzinski1992}
W.~Rudzinski and D.H.~Everett,
\textit{Adsorption of Gases on Heterogeneous Surfaces},
Academic Press, London (1992).

\bibitem{Klotz1971}
I.M.~Klotz and D.L.~Hunston,
``Properties of graphical representations of multiple classes of
binding sites,''
\textit{Biochemistry} \textbf{10}, 3065--3069 (1971).

\bibitem{Toth1971}
J.~T\'oth,
``State equations of the solid--gas interface layers,''
\textit{Acta Chim.\ Acad.\ Sci.\ Hung.}\ \textbf{69}, 311--328 (1971).

\bibitem{Jaroniec1975}
M.~Jaroniec,
``Adsorption on heterogeneous surfaces: the exponential equation
for the overall adsorption isotherm,''
\textit{Surface Sci.}\ \textbf{50}, 553--564 (1975).

\bibitem{Redlich1959}
O.~Redlich and D.L.~Peterson,
``A useful adsorption isotherm,''
\textit{J.\ Phys.\ Chem.}\ \textbf{63}, 1024 (1959).

\bibitem{BET1938}
S.~Brunauer, P.H.~Emmett, and E.~Teller,
``Adsorption of gases in multimolecular layers,''
\textit{J.\ Am.\ Chem.\ Soc.}\ \textbf{60}, 309--319 (1938).

\bibitem{Vieth1976}
W.R.~Vieth, J.M.~Howell, and J.H.~Hsieh,
``Dual sorption theory,''
\textit{J.\ Membr.\ Sci.}\ \textbf{1}, 177--220 (1976).

\bibitem{Nederlof1990}
M.M.~Nederlof, W.H.~van Riemsdijk, and L.K.~Koopal,
``Determination of adsorption affinity distributions,''
\textit{J.\ Colloid Interface Sci.}\ \textbf{135}, 410--426 (1990).

\bibitem{Cernik1995}
M.~Cernik, M.~Borkovec, and J.C.~Westall,
``Regularized least-squares methods for the calculation of discrete
and continuous affinity distributions for heterogeneous sorbents,''
\textit{Environ.\ Sci.\ Technol.}\ \textbf{29}, 413--425 (1995).

\end{thebibliography}

\end{document}
